% ============================================================
% Sección 5: Conclusiones
% ============================================================

\section{Conclusiones}
\label{sec:conclusiones}

Este trabajo presenta una implementación de Redes Neuronales Informadas por Física
con activación sinusoidal (PINN-SIREN) para la resolución de la ecuación de Helmholtz
2D con campo complejo, bajo un paradigma \textit{libre de malla} (\textit{mesh-free})
\cite{Karniadakis2021_review,Cuomo2022_review}.
Los dos experimentos completados (NB01 y NB02) validan la arquitectura con precisión
superior al estado del arte.

Las conclusiones principales son:

\begin{enumerate}
  \item \textbf{Precisión excepcional en 1D y 2D.}
    La arquitectura SIREN alcanzó errores $L^2$ de $0.006\%$ (1D) y $0.171\%$ (2D),
    varios órdenes de magnitud por debajo del umbral de tesis ($< 5\%$).

  \item \textbf{La estrategia Adam + L-BFGS es esencial.}
    Adam proporciona exploración global robusta; L-BFGS refina la solución con
    información de curvatura de segundo orden.
    El uso exclusivo de cualquiera de los dos optimizadores produce resultados inferiores.

  \item \textbf{El peso físico $\lambda = 0.1$ equilibra datos y física en 2D.}
    En el caso 2D, la función de pérdida física incluye dos residuos (real e imaginario),
    duplicando efectivamente el peso de la física si se usa $\lambda = 1.0$.
    La calibración a $\lambda = 0.1$ fue crítica para la convergencia.
\end{enumerate}

\subsection*{Trabajo futuro}

Las extensiones naturales de este trabajo incluyen:

\begin{itemize}
  \item \textbf{Simulación de speckle óptico} (NB03): aplicar el modelo validado en NB02
    a la generación de speckle, reemplazando la condición de onda plana por una frontera
    de fase aleatoria $\phi(x) \sim \mathcal{U}(0, 2\pi)$ en $y = 0$, y validar
    estadísticamente el patrón generado con el contraste de Goodman $|C - 1| < 0.1$.

  \item \textbf{Benchmark FEM} (NB04): comparación cuantitativa contra FEniCSx para
    múltiples realizaciones de speckle, cuantificando el \textit{Speed-up Factor}
    $S = T_{\mathrm{FEM}} / T_{\mathrm{PINN}}$ \cite{Thuerey2021_PBDL} y la
    huella de memoria pico (\textit{Peak Memory Footprint}), según los objetivos
    de la investigación.

  \item \textbf{Medios inhomogéneos} (NB05-NB06): extensión a la ecuación de Helmholtz
    con índice de refracción variable $n(\mathbf{r})$, relevante para speckle en
    materiales GRIN y tejidos biológicos.
\end{itemize}
